\section{Two Pell roots carry two quotients}
\label{sec:two-pell-carriers}

This section gives an independent compact improvement, retaining the exact
admissible942 coefficient tuple of the incoming \emph{Pell Carrier65}
package.  That package removes the first-index quotient by using the first
Pell root.  We now use a second Pell root to remove the native transport
quotient.  The factor four in an odd-ordinate congruence is what makes the
second reduction work without a parity restriction on the auxiliary base.
The resulting ordinary polynomial has thirteen strictly positive variables,
69 live binary arithmetic operations, and exact degree~174.  Its standard
unrestricted-integer presentation has~39 variables, 160 operations, and
degree~348.  A companion trades one extra operation for degrees170 and340.

The parent and its current calibration are retained under the compact
component's \texttt{inherited/} directory.  In particular the complete
admissible942 proof, rather than an arbitrary four-number tuple satisfying
elementary inequalities, gives these equations their consistency meaning.
The general use of a Pell ordinate as a quotient carrier already occurs
in package~170.  The present contribution is the stronger residue, its
use for native sign recovery before rank decoding, and its composition
with the first-root construction.

\subsection{The entire polynomial}

There are exactly thirteen independent positive integer inputs,
in the order used by the machine-readable equation:
\[
 J,F,\alpha,f,i,T,s,w,n,\eta,\zeta,m,\gamma>0.
\]
Retain the four fixed integers $\beta,K_0,M_C,M_F$ from the exact parent.
Their unsigned big-endian byte strings total~94,689 bytes, and their
largest bit length is~189,693.  The compact component includes their complete
finite reconstruction, its~399-character seed, and all four literal
byte strings.  The polynomial below contains no variable exponent,
division, floor, factorial, or unspecified predicate.  Every other letter
is a polynomial abbreviation:
\begin{align*}
 h&=\beta J,& q&=h+1,& C&=q-F-\alpha,\\
 X&=wq,&Y&=sq^3,&E&=XY,\\
 k&=\eta+\zeta,&\ell&=EkY,&c&=kY+\eta,\\
 a&=E+Y,&H&=4a+3,&\Delta&=a^2+H=(a+2)^2-1,\\
 D&=X+ac+\gamma H,&\tau&=En+k-R,&B&=(K_0+w)C-F,
\end{align*}
where
\[
 R=[h(q-F)+\alpha](q^2-1)+(M_C+qM_F)J.
\]
The last definitions are
\begin{align*}
 b&=cT-R,& V&=f^2b-c,&S&=hi\Delta c^2,\\
 L&=V^2,& y&=1+4\bigl[h(m-LB)-B\bigr],\\
 N_0&=\tau^2-\ell(\ell+k),&
 N_1&=D^2-\Delta c^2,\\
 N_a&=S^2(V^2-y^2)+y^2,&
 N_s&=\Delta f^2-S^2.
\end{align*}
The single equation is
\begin{equation}
 \boxed{P=N_0N_1N_aN_s-\Delta=0.}
 \label{eq:second-carrier}
\end{equation}
The companion changes only $L$ to $bV$.  Notice that $y$ is a signed
polynomial value, not an additional positive unknown.

\subsection{A residue with a factor of four}

For an integer $S\ge2$, write
\[
 \chi_S(r)+\psi_S(r)\sqrt{S^2-1}
   =(S+\sqrt{S^2-1})^r.
\]
The odd subsequence $u_j=\psi_S(2j+1)$ satisfies
\[
 u_0=1,\quad u_1=4S^2-1,\quad
 u_{j+2}=(4S^2-2)u_{j+1}-u_j.
\]
Induction modulo $4S^2$ therefore proves
\begin{equation}
 \psi_S(2j+1)\equiv(-1)^j\pmod{4S^2}.
 \label{eq:odd-four-residue}
\end{equation}
This holds for odd as well as even $S$.  The weaker congruence modulo
$S^2$ would be insufficient for the sign selection used below.

Suppose that integers $V,y$ satisfy
\begin{equation}
 S^2V^2-(S^2-1)y^2=1.\label{eq:second-adjacent}
\end{equation}
Neither coordinate is zero.  The positive solutions of
$x^2-(S^2-1)u^2=1$ are exactly
$(\chi_S(r),\psi_S(r))$, $r\ge1$.  One can prove this directly by
descent: for $u>1$, multiplication by $S-\sqrt{S^2-1}$ sends the
ordinate to $Su-x$, which is positive and smaller than $u$, because
$S-1<\sqrt{S^2-1}<x/u<S$.  The new abscissa is positive as well.
The terminal positive pair is $(S,1)$, and reversing the descent gives
the stated powers.

Apply this to $x=S|V|$ and $u=|y|$.  Modulo $S$, the even-index
$\chi$ values are $+1$ or $-1$, whereas the odd-index values are zero.
The index is consequently odd.  Equation\eqref{eq:odd-four-residue}
then gives
\[
 y\equiv+1\text{ or }-1\pmod{4S^2}.
\]
If $y\equiv1\pmod4$, the minus sign is impossible.  Thus
\begin{equation}
 y\equiv1\pmod4\quad\Longrightarrow\quad4S^2\mid y-1.
 \label{eq:carrier-full-divisibility}
\end{equation}
Also, directly from\eqref{eq:second-adjacent},
\begin{equation}
 |y|<2|V|.\label{eq:carrier-root-bound}
\end{equation}
No packed-rank interpretation or positivity of $V$ or $y$ was used.

\subsection{Soundness before rank recovery}

On the entire positive domain, $h\ge31$, $c>1$, $\Delta>1$ and
$S>c$.  At a zero of\eqref{eq:second-carrier}, cancellation of
the positive integer $\Delta$ gives
\[
 N_0N_1N_a\bigl[f^2-\Delta(hi c^2)^2\bigr]=1.
\]
All four factors are integer units.  Each negative sign is excluded
independently.  The consecutive-product negative-Pell obstruction excludes
$N_0=-1$, since
$\ell(\ell+k)=EY(EY+1)k^2$ and $EY\ge2$.  For completeness, if
$x^2-t(t+1)u^2=-1$ with $t\ge2$ and least positive $u$, then
$tu<x<(t+\tfrac12)u$ after taking absolute values; multiplication by
$(2t+1)-2\sqrt{t(t+1)}$ gives a positive smaller ordinate
$(2t+1)u-2x$, a contradiction.  This does not assume $\tau>0$.

Because $\Delta$ is zero or three modulo four, neither $N_1$ nor the
normalized strong norm can be $-1$.  Finally $N_a$ is congruent modulo
four to either $y^2$ or $V^2$, and cannot be $-1$.  Hence
\begin{equation}
 N_0=N_1=N_a=1,\qquad N_s=\Delta.
 \label{eq:four-unit-signs}
\end{equation}
In particular $N_a=1$ is\eqref{eq:second-adjacent}.  Our definition
of $y$ gives $y\equiv1\pmod4$, so
\eqref{eq:carrier-full-divisibility}, together with $h\mid S$, yields
\[
 h^2\mid h(m-LB)-B,\qquad\text{and therefore }h\mid B.
\]
No primality or parity assumption on $h$ enters this deduction.

For the main profile, $L=V^2\ge|V|\ge1$.  For the companion
$L=bV$, the same inequality holds at a zero.  Indeed $b=0$ would give
$V=-c$ and
\[
 y^2-c^2=\frac{c^2-1}{S^2-1}\in(0,1),
\]
impossible for integers.  Equation\eqref{eq:four-unit-signs} gives
$f^2=1+\Delta h^2i^2c^4>c$, so $V=f^2b-c$ has the sign of nonzero
$b$.  Therefore $bV=|b||V|\ge|V|$.

Set $W=1+hL>0$.  If $B\le-1$, positivity of $m$ forces
\[
 y=1+4(hm-WB)\ge1+4(h+W)>4hL\ge4|V|,
\]
contradicting\eqref{eq:carrier-root-bound}.  Thus $B\ge0$.
The omitted native quotient is consequently recovered as the positive
integer
\[
 z=B/h+1.
\]
Its parent factor is exactly
\[
 N_t=(K_0+w)C+(q-F)-zh=B+q-zh=1.
\]
Keep the other parent inputs, replace its strong multiplier by $hi$,
and replace its positive auxiliary ordinate by $|y|$.  Every parent
factor now has the required value, and the Carrier65 polynomial vanishes.
This is the complete soundness map.  It uses no conclusion about $C$,
$R$, or the first Pell index before reaching the already proved parent.

\subsection{Completeness with a fresh strong seed}

A finite syntactic ZFC refutation supplies the parent's canonical outer
data.  Among their properties are
\[
 N_0=N_1=1,\quad R\equiv3\pmod4,\quad
 c=\psi_{a+2}(R),\quad X=2^R,
\]
as well as $C>0$, $F<q$, $h\mid B$, and $R>3q$.  The inherited
first-root coordinate $n$ is positive.  These data are independent of
the private strong seed and the auxiliary coordinates.

There is actually a strict native surplus: $2^R>q^2$ gives $w=X/q>q$,
and hence
\begin{equation}
 B=(K_0+w)C-F>w-q>0.\label{eq:strict-native-surplus}
\end{equation}
Choose a positive solution
\[
 f^2-\Delta I^2=1,\qquad hc^2\mid I.
\]
Such solutions exist cofinally.  In the finite ring
$(\mathbb Z/(hc^2)\mathbb Z)[t]/(t^2-\Delta)$, the element $a+2+t$
is a unit of norm one and therefore has finite order.  Any positive
multiple of that order gives $I\equiv0\pmod{hc^2}$; its positive
powers give arbitrarily large choices.  Set $i=I/(hc^2)>0$ and
$S=\Delta I$.  Then $N_s=\Delta$ for our definitions.

Put
\[
 V_*=\frac{\chi_S(R)}{S},\qquad y_* =\psi_S(R).
\]
The quotient is a positive integer because $R$ is odd.  Write $R=2j+1$
and $\chi_S(R)/S=Q_j(S^2)$.  The elementary identities retained in
the parent proof are
\[
 Q_j(0)=(-1)^jR,\qquad
 Q_j(1-A^2)=(-1)^j\psi_A(R),\quad A=a+2.
\]
Since $j$ is odd, $S^2=1-A^2\pmod{f^2}$ and $c\mid S$, they give
$V_*\equiv-c\pmod{f^2}$ and $V_*\equiv-R\pmod c$.
Moreover $\gcd(c,f)=1$ and $f^2\equiv1\pmod c$.  Thus
\[
 T=\frac{V_*+c+Rf^2}{cf^2}>0
\]
is an integer and our polynomial value $V$ is $V_*$.  These quotients
construct witnesses; none is an operation in the delivered polynomial.

Equation\eqref{eq:odd-four-residue} gives
$y_*\equiv-1\pmod{4S^2}$.  Use the signed value $y=-y_*$, so
$y\equiv1\pmod{4S^2}$, and define
\begin{equation}
 m=\frac{y-1}{4h}+\frac{(1+hL)B}{h}.
 \label{eq:new-positive-carrier-witness}
\end{equation}
It is an integer.  It is positive because $B$ is a positive multiple of
$h$ by\eqref{eq:strict-native-surplus}, $L\ge V\ge1$, and
$|y|<2V$:
\[
 m> -\frac{2V+1}{4h}+1+hL>0.
\]
Equation\eqref{eq:new-positive-carrier-witness} makes the literal
carrier equal the chosen signed $y$, and all four factors in
\eqref{eq:four-unit-signs} have their required values.  Thus $P=0$.
This is an existence construction that refreshes the private Pell seed;
it is not a bijection on every old and new zero.

\subsection{Exact size and unrestricted integer variables}

Relative to the~65-operation parent, remove the three final native-factor
rows, add the one subtraction defining $B$, remove one product
multiplication, scale the strong multiplier by $h$ using one multiplication,
and compute $y$ with six operations.  The count is therefore
$65-3+1-1+1+6=69$.  The companion requires one extra multiplication for
$bV$.  Squares and scalar multiplications are charged as binary operations.

The respective factor degrees are
\[
\begin{array}{c|rrrr|r}
 &N_0&N_1&N_a&N_s&P\\\hline
L=V^2&22&18&86&48&174\\
L=bV&22&18&82&48&170
\end{array}
\]
The main-norm degree uses the exact cancellation, with $Z=X+\gamma H$,
\[
 (ac+Z)^2-(a^2+H)c^2=2acZ+Z^2-Hc^2.
\]
Let $q_0=\beta J$ and $k_0=\eta+\zeta$.  The complete leading
homogeneous form of the main equation is
\[
 -128q_0^{99}\gamma i^4k_0^{15}w^{14}s^{27}f^8T^4
          (q_0-F-\alpha)^2.
\]
The companion replaces $f^8$ by $f^4$.  The coefficient selecting
$\eta^{15}$ and $F^2$ is $-128\beta^{99}$, which is nonzero on the
actual fixed tuple.  These are exact degrees, not only graph estimates.

To obtain unrestricted integer equations, replace each positive input by
\[
 1+u^2+v^2+w(w+1),\qquad u,v,w\in\mathbb Z.
\]
The classical three-square theorem applied to $4N+1$ proves that every
$N\ge0$ is $u^2+v^2+w(w+1)$: exactly one of the three squares is odd,
and subtracting its1 and dividing by4 gives the representation.
Conversely this expression is nonnegative on all integer triples.
The substitution uses seven operations per coordinate.  It gives
$39$ integer variables and respectively $160$ or $161$ operations.
Degrees double to348 and340.  Nonzero highest homogeneous terms survive:
setting each second and third new coordinate to zero doubles the
monomial exponents injectively.

The independent compact checker verifies~25 or~26 exact formula cuts,
the full~96-term leading forms, all coordinate and coefficient bindings,
the actual fixed-tuple degree witnesses, four coefficient reconstructions,
and~1,278 finite Pell-carrier fixtures.  Targeted changed gates are rejected.
These finite checks complement the preceding unbounded proof.  They do
not replace the inherited native/K215/refutation calibration or provide a
fresh end-to-end proof-assistant derivation.  Composing that calibration
with the two witness maps and the exact domain translation proves in ZFC
that the integer equation has no zero if and only if
$\operatorname{Con}(\mathrm{ZFC})$.

The result is a tradeoff against the matched parent and inspected compact
family.  Package~170 has a~13-positive-variable, 70-operation degree~206
point and a~71-operation degree~162 point, with a different~295,388-byte
tuple.  The present~69-operation degree~174 point does not assert a
minimum over all Diophantine equations, all integer domains, or all
coefficient architectures.
